\section{Methodology}
\label{s:methodology}

The analysis measures the impact of the Brasília Subway System on local economic outcomes at census-tract level. The main outcome is income growth, defined as the change in (log) per-capita income between 2000 and 2010. Let $i$ index census tracts and let $r(i)$ denote the administrative region (Região Administrativa, RA) to which tract $i$ belongs. The dependent variable is constructed as a cross-section of changes\footnote{A ``cross-section of changes'' means that each tract appears once in the regression, with its outcome already expressed as the 2000-2010 change. This differs from a panel specification (which would stack observations for 2000 and 2010 and include time variation explicitly). In our specification, the time dimension is absorbed into the construction of $\Delta y_{i}$.}:
\begin{equation}
    \Delta y_{i} \equiv \log (\text{income}_{i,2010}) - \log (\text{income}_{i,2000})
    \label{eq:incomechange}
\end{equation}

Let $D_{i}$ denote exposure to the subway. Subway exposure is defined using geographic proximity to stations. For each census tract $i$, let $d_{i}$ denote the Euclidean distance (in meters) from the tract (its centroid) to the nearest subway station open at the time of the 2010 Census. The baseline exposure indicator classifies a tract as treated if it lies within 1,000 meters of the nearest station:
\begin{equation}
    D_{i} \equiv 1\{d_{i} \leq 1000\}
    \label{eq:exposure_var}
\end{equation}
thus a tract is classified as exposed ($D_{i} = 1$) if it lies within 1,000 meters of the nearest station, and unexposed ($D_{i} = 0$) otherwise.\footnote{This threshold follows the Institute for Transportation and Development Policy (ITDP) \emph{People Near Transit} metric and corresponds roughly to a 10--15 minute walk, as discussed in Section~\ref{s:data}.} 

A central concern is that subway placement is endogenous: lines and stations may have been located in areas with different pre-trends, political priorities, or expected growth. To address this, the empirical strategy uses two-stage least squares (2SLS). The instrument $Z_{i}$ is based on a discarded/planned subway project (the ``counterfactual'' design), capturing intended network placement that was not fully implemented.\footnote{The planned-project instrument is constructed as a binary proximity measure, analogous to the realized exposure. Because planned station locations are not available, planned access is derived from the minimum distance between the census-tract centroid and the closest segment of the planned (discarded) subway alignment. The instrument equals one if this distance is below 1,000 meters and zero otherwise, matching the treatment threshold. This represents a departure from the realized exposure variable, which is defined using distance to the nearest station rather than to the nearest line segment. The asymmetry introduces classical measurement error into the instrument: $Z_{i}$ is a noisy proxy for the latent concept of planned station proximity, which attenuates the first-stage coefficient and biases the 2SLS estimates toward zero. The estimated income effects should therefore be interpreted as a lower bound on the true causal effect.} The exclusion restriction requires that proximity to the IMT alignment affects income growth only through realized subway access, and not through any independent channel. The defense of this assumption rests on three arguments.

The first argument is temporal. The IMT alignment was designed in 1986--1987, more than a decade before the baseline year of this study (2000) and nearly a quarter century before the end of the outcome period (2010). Income growth trajectories at the census-tract level between 2000 and 2010 could not have influenced a planning document drafted in the mid-1980s: reverse causality is ruled out by construction. This approach follows an established strategy in the empirical literature on urban transportation \citep{baumsnow2007highways, severen2023commuting, mayertrevien2017largescale, brinkman2024freeway}: network plans drawn to address the conditions of their era are used to isolate the causal effects of infrastructure on later outcomes, precisely because the fundamentals that shaped the plan---population distribution, transport connectivity, and topography as they existed at the time of planning---predate the outcome dynamics of interest.

The second argument concerns the reasons why the IMT plan was never implemented. A plan that predates the outcome still needs to be abandoned for reasons unrelated to subsequent income growth. In the Brasília case, the available evidence points to three institutional factors: a change of state administration between the completion of the IMT study in early 1987 and the eventual construction phase that began in 1992; the severe financing constraints of the period, marked by hyperinflation and contested federal transfers; and a technical reassessment of the system's mode and alignment by the incoming administration. None of these factors is plausibly linked to the income trajectories of specific census tracts over the 2000--2010 decade. This mirrors the pattern documented in the closest analogues in the literature: \cite{mayertrevien2017largescale} show that the Paris RER plan diverged from its initial design following a political transition and budgetary pressures, and \cite{brinkman2024freeway} document that U.S. freeway construction departed from initial plans due to organized community opposition and subsequent regulatory changes---in both cases, the divergence reflects factors that are unrelated to the neighborhood-level outcomes being estimated.

The third argument addresses the geographic overlap between the IMT alignment and the realized network. Both proposals were designed to connect the \textit{Plano Piloto} to the same peripheral administrative regions, so some similarity in their general corridor is expected and does not invalidate the instrument. What matters for identification is not that the two alignments are distinct at the corridor level, but that the specific tracts classified as treated differ between them, and that these differences arise from routing decisions unrelated to income growth.\footnote{The data confirm this. Among the 1,887 tracts within 10~km of a station, the realized treatment and the planned instrument disagree for 16.2\% of tracts: 21\% of treated tracts fall outside the instrument and 28\% of instrument-exposed tracts are untreated. The two indicators are correlated at 0.64, leaving ample independent variation, which the strong first stage confirms is informative (the bivariate first-stage $F$-statistic is about 1{,}284; with administrative-region fixed effects and controls, the first-stage $F$-statistics in Table~\ref{tab:iv_results} range from 733 to 816).} A tract may fall within 1,000 meters of the IMT alignment but outside that range from the realized network, or the reverse, depending on the local routing choices of each design. Those differences---a curve adjusted for topography, a segment rerouted after a change in technical criteria---reflect the administrative and engineering decisions of 1986--1987, not the income trajectory of the tracts involved.

The IV strategy thus relies on (i) instrument relevance, i.e., planned exposure $Z_{i}^{(k)}$ must predict realized exposure $D_{i}^{(k)}$; this is assessed using the first-stage F-statistic; and (ii) the exclusion restriction, namely that---conditional on baseline covariates $X'_{i,2000}$ and administrative-region fixed effects $\alpha_{r(i)}$---planned exposure affects income growth only through realized subway exposure. The 2SLS coefficient $\beta_{1}$ is interpreted as a local average treatment effect (LATE) for census tracts whose realized exposure is shifted by the planned alignment. As we discussed in Section~\ref{s:data}, the preferred specification restricts the sample to tracts within 10 km of the nearest subway station to improve comparability among nearby neighborhoods.

The estimation is carried out via 2SLS with administrative-region (RA) fixed effects. The first-stage equation is:
\begin{align}
    D_{i}^{(k)} &= \pi_{0} + \pi_{1}Z_{i}^{(k)} + X'_{i,2000}\vartheta + \alpha_{r(i)} + u_{i}, \label{eq:first_stage_rafe}
\end{align}
where $D_{i}^{(k)} \in \{0,1\}$ is the binary subway exposure indicator that equals one if tract $i$ lies within $k$ meters of the nearest station and zero otherwise; $Z_{i}^{(k)} \in \{0,1\}$ is the corresponding binary instrument constructed from the planned/discarded network (defined analogously for the same threshold $k$); $X'_{i,2000}$ is the vector of baseline covariates (log distance to Brasília's center in 2000, log distance to the nearest highway in 2000, share of household heads with completed higher education in 2000, illiteracy share in 2000, share aged 65$+$ in 2000, log per-capita income in 2000 and log population in 2000); and $\alpha_{r(i)}$ are RA fixed effects. The fitted values from \eqref{eq:first_stage_rafe}, $\hat{D}_{i}^{(k)}$, capture the component of realized subway exposure explained by the planned/discarded network, conditional on baseline covariates and RA-level unobservables.

The second-stage equation is:
\begin{align}
    \Delta y_{i} &= \beta_{0} + \beta_{1} \hat{D}_{i}^{(k)} + X'_{i,2000}\gamma + \alpha_{r(i)} + \varepsilon_{i}, \label{eq:second_stage_rafe}
\end{align}
where $\Delta y_{i}$ is income growth from 2000 to 2010. This equation relates medium-run income growth across census tracts to \emph{predicted} subway exposure, purged of endogenous placement decisions through the instrumental-variable strategy. The coefficient of interest, $\beta_{1}$, measures the causal effect of being exposed to the subway system—within distance threshold $k$—on income growth over the 2000--2010 decade, using only variation in exposure induced by the planned/discarded network. As a result, $\beta_{1}$ captures how income growth differs between tracts that became effectively exposed to the subway and otherwise similar tracts that did not, holding constant baseline socioeconomic characteristics and unobserved RA-level factors. In other words, the estimation assesses whether census tracts exposed to the subway experienced a larger (or smaller) positive change in income over the decade relative to non-exposed tracts.

Administrative-region fixed effects are included to absorb unobserved factors that are common to all census tracts within the same Região Administrativa (RA) and that may systematically affect income growth between 2000 and 2010. Even though the dependent variable is a change $\Delta y_{i}$, different RAs can follow different average growth paths over the decade due to persistent differences in local public services, land-use regulation, infrastructure provision, labor-market conditions, and broader neighborhood amenities. By adding RA fixed effects, identification comes from comparing more and less subway-exposed tracts within the same RA, reducing bias from cross-RA differences that could otherwise be correlated with subway placement or with the planned-network instrument. Year fixed effects are not included, as the specification is estimated in long differences.\footnote{A long-differences specification collapses the time dimension into the outcome variable, expressing the dependent variable as the change between two periods. With a single cross-section of changes, a year fixed effect would be mechanically collinear and therefore redundant.}

Standard errors are clustered at the RA level. Clustering at the RA level allows for arbitrary correlation in unobserved factors within administrative regions, which is appropriate given that tracts in the same RA share common policies and local conditions.
